\scanpage{171}\mspage{116}
\[
f^*:\mathrm{HOT}_W(S)\longrightarrow\mathrm{HOT}_W(S')\tag*{$(*)$}
\]
et le foncteur induit
\[
f^*:\mathrm{HOT}^{\mathrm{lc}}_W(S)\unsure{\longrightarrow|\text{---}}\mathrm{HOT}^{\mathrm{lc}}_W(S').
\]

Quand on décrit les catégories $\mathrm{HOT}(S)$, conformément à la recette générale, à l'aide de \inserted{\unsure{contrav.}} foncteurs
\[
F:S^{\unsure{o|0|\circ}}\longrightarrow\mathrm{Cat}
\]
(préfaisceaux en catégories sur $S$), ou ce qui revient au $=$, de catégories scindées, la description de la \deleted{première} flèche $(*)$ est évidente : elle revient à la loi
\[
\left|\begin{aligned}
&F\longmapsto F\circ f^o\\
&\underline{\mathrm{Hom}}(S^o,\mathrm{Cat})\unsure{\text{---}|\longrightarrow}\underline{\mathrm{Hom}}(S'{}^o,\mathrm{Cat}),
\end{aligned}\right.
\]
en passant aux localisés \unsure{pour la|pour les} $W$-\unsure{équivalences|équivalence} argument par argument. Dans le langage des catégories fibrées, il est bien connu que cette opération triviale sur les préfaisceaux en \unsure{catégorie|catégories}, correspond à l'opération \uline{image inverse} de catégories sur $S$ -- laquelle opération passe aux scindages.
Cela nous donne \deleted{\unsure{donc}}
\[
f^*(\mathrm{hot}_{\deleted{\unsure{X}}W,S}(X))=\mathrm{hot}_{W,S'}(\underbrace{X\times_S S'}_{X'}),
\]
pour \deleted{\unsure{tout}} $X$ fibrée scindée sur $S$. Or on sait que le scindage \inserted{de $X$} ne joue aucun rôle, \unsure{finalement|justement}, dans la description de la classe
\[
\mathrm{hot}_{W,S}(X)\in\mathrm{HOT}_W(S)
\]
% archive sheet number: 170
